% Figure 0.1 : les trois semestres comme une suite de problèmes résolus
\documentclass[border=6pt]{standalone}
\usepackage{coursfig}
\begin{document}
\begin{tikzpicture}[node distance=5mm and 9mm]
  \tikzset{blk/.style={box=#1, minimum width=52mm, minimum height=15mm, align=center}}
  % colonne = semestre ; chaque bloc : outil (gras) + problème vécu (italique)
  \newcommand{\blkt}[2]{\textbf{#1}\\[-.3mm]{\normalsize\itshape\color{Muted}#2}}

  \node[blk=Ocre] (a) {\blkt{Déploiement manuel}{un serveur, tout à la main}};
  \node[blk=Ocre, below=of a] (b) {\blkt{Multi-machines}{la configuration diverge}};
  \node[blk=Ocre, below=of b] (c) {\blkt{Infrastructure as Code}{Vagrant, Ansible, Terraform}};

  \node[blk=Enc, right=of a] (d) {\blkt{Conteneurisation}{Podman, images OCI}};
  \node[blk=Enc, below=of d] (e) {\blkt{Orchestration}{Kubernetes, réconciliation}};
  \node[blk=Enc, below=of e] (f) {\blkt{CI/CD, GitOps}{observabilité}};

  \node[blk=Sea, right=of d] (g) {\blkt{Reproductibilité}{données et modèles, DVC}};
  \node[blk=Sea, below=of g] (h) {\blkt{MLflow, Airflow}{pipelines d'entraînement}};
  \node[blk=Sea, below=of h] (i) {\blkt{Monitoring de dérive}{Spark, Kafka}};

  \foreach \u/\v in {a/b,b/c,d/e,e/f,g/h,h/i} \draw[flow=Muted] (\u) -- (\v);
  % passage d'un semestre au suivant : du bas d'une colonne au haut de la suivante
  \draw[flow=Ink, rounded corners=4pt] (c.east) -- ++(4.5mm,0) |- (d.west);
  \draw[flow=Ink, rounded corners=4pt] (f.east) -- ++(4.5mm,0) |- (g.west);

  \foreach \n/\t/\c/\x in {a/SEMESTRE 1 \textperiodcentered{} FONDATIONS/Ocre/a,
                            d/{SEMESTRE 2 \textperiodcentered{} CONTENEURS ET CI\slash CD}/Enc/d,
                            g/SEMESTRE 3 \textperiodcentered{} MLOPS ET BIG DATA/Sea/g}
     \node[ftitle, text=\c, above=2mm of \x] {\t};
\end{tikzpicture}
\end{document}
